\chapter{\stateoftheartname}
\label{chap:marco-ireb}

\section{Ingenier\'ia de requisitos basada en normas}

\subsection{IREB y el CPRE Foundation Level Handbook}

El International Requirements Engineering Board (IREB) define un marco de referencia para la ingenier\'ia de requisitos a trav\'es de su programa de certificaci\'on CPRE. El \emph{Foundation Level Handbook}~\cite{cpre} establece los fundamentos conceptuales y las buenas pr\'acticas de la disciplina.

El handbook define \textbf{requisito} desde tres puntos de vista complementarios: como una necesidad percibida por un \emph{stakeholder}, como una capacidad o propiedad que el sistema debe tener, y como la representaci\'on documentada de esa necesidad o capacidad~\cite[p.~10]{cpre}. Define \textbf{Requirements Engineering} como el enfoque sistem\'atico y disciplinado para especificar y gestionar requisitos con el objetivo de entender los deseos y necesidades de los \emph{stakeholders} y reducir el riesgo de entregar un sistema inadecuado~\cite[pp.~10--11]{cpre}.

\subsubsection{Las tres actividades fundamentales}

IREB distingue tres actividades fundamentales~\cite[pp.~16--28]{cpre}:

\begin{description}
    \item[Elicitaci\'on:] el proceso de buscar, capturar y consolidar requisitos desde fuentes disponibles. Las fuentes t\'ipicas incluyen \emph{stakeholders}, documentos, sistemas existentes y observaciones.
    \item[Documentaci\'on:] la especificaci\'on estructurada de requisitos siguiendo criterios de calidad. IREB recomienda el uso de plantillas y atributos (ID, nombre, fuente, prioridad, \emph{rationale}, estado, versi\'on).
    \item[Validaci\'on:] el proceso de confirmar que los requisitos documentados reflejan las necesidades reales de los \emph{stakeholders}.
\end{description}

A estas tres actividades se a\~nade una transversal de \textbf{gesti\'on de requisitos}, que abarca almacenamiento, versionado, trazabilidad y control de cambios.

\subsubsection{Los nueve principios fundamentales}

El handbook establece nueve principios fundamentales~\cite[pp.~16--18]{cpre}: (1) orientaci\'on al valor, (2) implicaci\'on de \emph{stakeholders}, (3) entendimiento compartido, (4) contexto del sistema, (5) problema-requisito-soluci\'on, (6) validaci\'on, (7) evoluci\'on, (8) innovaci\'on, y (9) trabajo sistem\'atico y disciplinado.

\subsection{ISO/IEC/IEEE 29148: criterios de calidad}

La norma ISO/IEC/IEEE 29148~\cite{iso29148} define los criterios que debe cumplir un requisito bien formado: necesario, singular, completo, verificable, no ambiguo, consistente y trazable.

\subsection{INCOSE Guide for Writing Requirements}

La INCOSE Guide~\cite{incose} recomienda la estructura ``El sistema deber\'a [verbo] [objeto] [condici\'on]'', verbos observables, ausencia de t\'erminos subjetivos y autosuficiencia del requisito.

\section{SpecKit y Spec-Driven Development}

\textsc{SpecKit} es un \emph{toolkit} de c\'odigo abierto (MIT) desarrollado por GitHub para practicar SDD. Su repositorio oficial cuenta con m\'as de 112\,000 estrellas, 229 contribuyentes y 162 \emph{releases} (v0.10.2, junio 2026). Est\'a escrito principalmente en Python (94.7\,\%) con scripts auxiliares en Shell y PowerShell~\cite{speckit_repo}.

La filosof\'ia SDD se fundamenta en seis principios~\cite{speckit_sdd}: (1) especificaciones como lengua franca, (2) especificaciones ejecutables, (3) refinamiento continuo, (4) contexto impulsado por investigaci\'on, (5) retroalimentaci\'on bidireccional y (6) ramificaci\'on para exploraci\'on. El pipeline b\'asico consta de cinco comandos: \texttt{constitution}, \texttt{specify}, \texttt{plan}, \texttt{tasks} e \texttt{implement}, a los que se a\~naden tres opcionales de calidad: \texttt{clarify}, \texttt{checklist} y \texttt{analyze}~\cite{speckit_docs}.

Un aspecto diferencial de \textsc{SpecKit} es su arquitectura extensible, que permite personalizar el flujo mediante \textbf{presets} (22 comunitarios, sobreescriben plantillas), \textbf{extensions} (105 comunitarias, a\~naden comandos como integraci\'on Jira o compuertas de calidad) y \textbf{workflows} (pipeline automatizado en YAML con compuertas de revisi\'on humana)~\cite{speckit_repo}.

\section{Herramientas profesionales de gesti\'on de requisitos}

Herramientas como IBM DOORS, Jama Connect o Siemens Polarion son el est\'andar industrial para la gesti\'on de requisitos. La tabla~\ref{tab:profesionales} compara sus capacidades con las de \textsc{SpecKit}.

\begin{table}[h]
\centering
\caption{Comparativa: herramientas profesionales vs. SpecKit}
\label{tab:profesionales}
\begin{tabularx}{\textwidth}{lcccc}
\toprule
\textbf{Caracter\'istica} & \textbf{DOORS} & \textbf{Jama} & \textbf{Polarion} & \textbf{SpecKit} \\
\midrule
Atributos personalizables & \checkmark & \checkmark & \checkmark & $\bigcirc$ \\
Clasificaci\'on por tipo & \checkmark & \checkmark & \checkmark & $\bigcirc$ \\
Trazabilidad bidireccional & \checkmark & \checkmark & \checkmark & $\bigcirc$ \\
Baselines & \checkmark & \checkmark & \checkmark & $\bigcirc$ \\
Workflow aprobaci\'on & \checkmark & \checkmark & \checkmark & $\bullet$ \\
Generaci\'on de c\'odigo & $\bigcirc$ & $\bigcirc$ & $\bigcirc$ & \checkmark \\
Integraci\'on IA nativa & $\bigcirc$ & $\bigcirc$ & $\bullet$ & \checkmark \\
\bottomrule
\end{tabularx}
\end{table}

Ninguna herramienta profesional integra generaci\'on de c\'odigo asistida por IA, que es precisamente el valor diferencial de \textsc{SpecKit}. Esta complementariedad entre gesti\'on y generaci\'on motiva el presente mapeo.

\section{Trabajos relacionados}

El uso de proyectos \emph{Open Source} como corpus para estudios emp\'iricos de ingenier\'ia del software es una pr\'actica consolidada en la literatura (conferencias MSR, miner\'ia de repositorios). Sin embargo, no existe hasta la fecha un estudio que mapee sistem\'aticamente un pipeline SDD contra un marco RE normativo utilizando requisitos derivados de fuentes documentales reales. Este trabajo cubre ese vac\'io.